\section{Related Academic Work}
\label{sec:related-work}
This section reviews the work most directly relevant to the thesis along three threads: the prior projects within the Reading the Reader initiative, the use of eye-tracking as a real-time signal in reading research, and the design and evaluation of typographic interventions.
Rather than cataloguing each source in isolation, the discussion draws out where the findings converge, namely that adaptive reading is feasible and that intervention design measurably shapes the reading experience, and where they leave a gap that this thesis addresses.

\subsection{Prior Work in the Reading the Reader Project}
\label{subsec:related-rtr-prior-work}
This thesis continues a line of student projects within the Reading the Reader initiative, and its contribution is best understood relative to that lineage.
Refsgaard and Farooq built the first flexible reading application and gaze data-collection platform for the project, conducting an experiment with 23 participants to assess the feasibility of the interface and data pipeline \cite{refsgaard2024rtr}.
Building directly on that codebase, Pereira and Andrei developed an adaptive e-reader that triggers typographic adjustments and evaluated four intervention designs (an Undo, a Popup, a Notification, and a Gradual change) through an eye-tracking study, reporting that careful typographic changes can improve reading speed and comfort while differing markedly in disruption and user acceptance \cite{pereira2024typography}.
In parallel, Palinec developed a reactive adaptive frontend intended to consume dynamic machine-learning output and adjust text presentation in real time \cite{palinec2024reactive}.

These projects share a common limitation that motivates the present work.
Each delivered a largely self-contained prototype, and the work evolved incrementally, with each project extending the previous frontend rather than building against a stable platform contract; Pereira and Andrei, for example, based their prototype on the earlier Refsgaard and Farooq code \cite{refsgaard2024rtr, pereira2024typography}.
This thesis responds by separating sensing, analysis, decision strategy, and intervention execution into independently replaceable modules, so that future intervention designs and decision providers can be integrated without rewriting the reading runtime or the researcher workflow.

\subsection{Eye-Tracking in Reading Research}
\label{subsec:related-eyetracking-reading}
A body of work establishes that eye movements are a productive signal for understanding and enhancing reading.
The eyeBook demonstrated an interactive reading application that detects the reader's position on the screen from gaze and generates contextual feedback, illustrating that gaze can drive in-situ responses during reading rather than only post-hoc analysis \cite{biedert2010eyebook}.
Subsequent work has shown that eye-movement behaviour can predict reading performance in online contexts \cite{guan2022predictive} and that reading versus not-reading behaviour can be classified automatically from eye-movement features \cite{landsmann2019classification}.
These results depend on first reducing raw gaze to oculomotor events using identification algorithms of the kind catalogued by Salvucci and Goldberg \cite{salvucci2000identifying}.
Collectively they support the architectural premise of this thesis: that a real-time analysis layer converting gaze into fixations, saccades, and regressions is a necessary foundation for any decision strategy that acts on reading behaviour.

\subsection{Context Preservation and Intervention Design}
\label{subsec:related-context-preservation}
The ETRA 2025 paper on context preservation formalises a key HCI challenge: adaptive reading interfaces can either assist or hinder reading depending on how changes are introduced and how recovery is supported.
In a within-subjects study with 22 participants, Jensen et al. report differences in RRT across intervention types \cite{ref2}. % \todo{confirm whether the paper reports these as statistically significant, and add the test/p-value if so}
The design of the interventions themselves has also been examined within the Reading the Reader project: Pereira and Andrei compared four intervention forms (Undo, Popup, Notification, and Gradual) and reported that they differ markedly in disruption and user acceptance \cite{pereira2024typography}.
The typographic parameters these interventions manipulate are in turn grounded in controlled readability studies, which show that font size and line spacing \cite{rello2016makeitbig} and font size and type \cite{beymer2008eyetracking} measurably affect online reading.

This supports the thesis argument that intervention design is a measurable system behaviour that must be architecturally supported through stable integration hooks and logging mechanisms.

\subsection{Transparent and Reproducible Analysis Pipelines}
\label{subsec:related-reading-analytics}
A theme common to the studies above is that the value of gaze analysis rests not only on which metrics are computed but on how reproducibly they are produced.
Each result depends first on reducing raw gaze to oculomotor events through an explicit identification algorithm, for which Salvucci and Goldberg provide a widely used taxonomy organised by the criterion each applies, such as velocity (I-VT) or dispersion (I-DT) \cite{salvucci2000identifying}.
Only once this extraction step is fixed and documented can downstream results, such as predicting reading performance \cite{guan2022predictive} or classifying reading versus non-reading behaviour \cite{landsmann2019classification}, be meaningfully compared or replicated across studies.

This carries a direct architectural consequence for the thesis: the analysis layer must be an explicit, inspectable stage rather than logic buried inside the reading runtime, so that the events driving every intervention can be logged, audited, and replayed.
A modular and transparent pipeline is therefore not only good experimental practice but a precondition for the systematic comparison of interventions and decision strategies that motivates this work.
